\begin{abstract}
Direct parameter-space merging of independently fine-tuned models has emerged as a computationally efficient paradigm for multi-task learning. However, during test-time adaptation (TTA), where merging coefficients are optimized online via entropy minimization to adapt to local data streams, the optimization trajectory is highly susceptible to transductive noise, leading to high-frequency spatial oscillations across network depth. We refer to this phenomenon as the \emph{Overfitting-Optimizer Paradox}: unsupervised local optimization fits transductive noise, causing catastrophic representation collapse and severe performance degradation. To resolve this paradox, we introduce \textbf{Riemannian Curvature-Regularized Test-Time Model Merging (RCR-Merge)}, a novel second-order geometric framework. We model the deep neural network parameter space as a Riemannian manifold where distance is locally scaled by the trace of the pre-trained base model's Fisher Information Matrix (FIM). During test-time adaptation, RCR-Merge regularizes the spatial total variation of merging coefficients across layer depth, dynamically scaling the penalty on adjacent layers by the geometric mean of their pre-trained base curvatures. We prove theoretically that this curvature-guided penalty acts as an analytical barrier protecting highly sensitive layers from representation-destroying coefficient fluctuations. Empirically, we evaluate RCR-Merge on a rigorous 30-seed simulation study over a synthetic Coupled Model II Landscape emulator that models multi-task adaptation across MNIST, FashionMNIST, CIFAR-10, and SVHN task profiles. RCR-Merge resolves the overfitting collapse, outperforming the static Uniform Baseline by \textbf{+3.06\% absolute}, unconstrained AdaMerging by \textbf{+10.33\% absolute}, and flat total variation regularization by \textbf{+5.01\% absolute}, while dramatically reducing optimization variance and stabilizing the adaptation trajectory. Our work establishes second-order landscape geometry as a vital protector of representation integrity during adaptive model merging.
\end{abstract}